\documentclass[11pt,a4paper]{extarticle}

\usepackage{kotex}
\usepackage[margin=17mm]{geometry}
\usepackage{amsmath, amssymb}
\usepackage{array,booktabs,tabularx}
\usepackage{enumitem}
\usepackage[table]{xcolor}
\usepackage[hidelinks]{hyperref}

\setlength{\parindent}{0pt}
\setlength{\parskip}{4pt}
\setlength{\abovedisplayskip}{5pt}
\setlength{\belowdisplayskip}{5pt}
\setlength{\abovedisplayshortskip}{4pt}
\setlength{\belowdisplayshortskip}{4pt}
\setlist[itemize]{leftmargin=18pt,itemsep=2pt,topsep=2pt}
\setcounter{tocdepth}{1}

\definecolor{lightgray}{RGB}{246,247,248}

\title{zkDPP 연구 맥락과 설계의 발전 과정}
\author{Research Context and Design Evolution}
\date{2026-07-24}
\renewcommand{\contentsname}{목차}

\begin{document}
\maketitle

\section*{문서 목적}

zkDPP 연구에서 고민한 문제와 설계 방향을 기록한다. 1, 2장은 선택적으로
읽는 배경이며, 현재 연구를 이해하기 위한 필수 흐름은 아니다.

\tableofcontents
\newpage

\section{연구의 출발}

\subsection{최초 목표}

최초 목표는 \textbf{원자재를 거래하는 플랫폼}을 만드는 것이었다. 이 플랫폼의
이름은 \textbf{zkDPP}로 정했었다. 거래 대상은 원자재부터 진행하여
모든 제조 공정을 마친 완제품까지 확장하는 구조로 생각했다.

\begin{itemize}
  \item 원자재의 정보를 투명하게 관리하고 거래한다.
  \item 원자재의 정보는 제조 공정을 거치며 누적되고 변경된다.
  \item 이러한 정보의 흐름은 완제품이 만들어질 때까지 투명하게 관리된다.
  \item 이처럼 누적되고 변경되는 제품 정보 자체를 DPP로 본다.
\end{itemize}

\subsection{플랫폼 구조}

\begin{center}
  \renewcommand{\arraystretch}{1.25}
  \begin{tabularx}{\textwidth}{@{}lX@{}}
    \toprule
    \textbf{객체}         & \textbf{역할}                       \\
    \midrule
    \textbf{Issuer}     & 사용자를 확인하고 Credential을 발급한다.       \\
    \textbf{사용자}        & 인증된 참여자이며, 거래에 따라 판매자 또는 구매자가 된다. \\
    \textbf{Server}     & 제품 정보, 생산 정보, 상품 검색과 구매 요청을 관리한다. \\
    \textbf{Blockchain} & 사용자와 각 사용자가 보유한 물품 목록을 기록한다.      \\
    \bottomrule
  \end{tabularx}
\end{center}

\subsection{다섯 가지 이벤트}

\[
  \boxed{\text{사용자 등록}}
  \;\longrightarrow\;
  \boxed{\text{제품 생산}}
  \;\longrightarrow\;
  \boxed{\text{제품 등록}}
  \;\longrightarrow\;
  \boxed{\text{제품 구매}}
  \;\longrightarrow\;
  \boxed{\text{제품 수령}}
\]

\textbf{1. 사용자 등록}

\[
  \begin{gathered}
    \text{Issuer의 사용자 확인 및 Credential 발급}
    \longrightarrow\;
    \text{Blockchain 및 Server에 사용자 등록}
  \end{gathered}
\]

사용자를 Server와 Blockchain에 등록하기 위한 요청이다.

Server 및 Blockchain에는 등록된 사용자의 물품 목록을 관리할 수 있는 영역을 생성한다.

\textbf{2. 제품 생산}

\[
  \begin{gathered}
    \text{제품 생산}
    \longrightarrow\;
    \text{사용자가 자신의 DB에 제품 생산 정보 기록}
    \;\longrightarrow\;
    \text{총 제품 개수 등 DB 정보 갱신}
  \end{gathered}
\]

\textbf{3. 제품 등록}

\[
  \text{판매자의 제품 생산}
  \;\longrightarrow\;
  \text{Blockchain의 판매자 물품 목록에 등록}
  \;\longrightarrow\;
  \text{Server에 동일한 제품 등록}
\]

\textbf{4. 제품 구매}

\[
  \begin{gathered}
    \text{구매자가 Server에서 상품 확인}
    \;\longrightarrow\;
    \text{판매자에게 구매 요청}
    \longrightarrow \\
    \text{판매자가 Blockchain에 자신의 물품 목록 갱신 요청}
    \longrightarrow \\
    \text{판매자가 거래용 }cm,nf\text{ 준비}
    \;\longrightarrow\;
    \text{구매자에게 전달}
  \end{gathered}
\]

\textbf{5. 제품 수령}

\[
  \begin{gathered}
    \text{구매자가 실제 제품 수령}
    \;\longrightarrow\;
    \text{구매자가 Blockchain에 자신의 물품 목록 갱신 요청} \\
    \;\longrightarrow\;
    \text{구매자의 보유 상태에 반영}
  \end{gathered}
\]

\section{상태 전이 모델}

zkDPP 플랫폼을 구현하기 위한 모델로 제품의 Lifecycle을 State Machine으로
표현하고자 했다. 이 모델에서, 암호학적인 기여가 추가적으로 필요하다고 보았다.

\[
  \boxed{\text{제품 Lifecycle}}
  \;\longrightarrow\;
  \boxed{\text{Lifecycle State Machine}}
\]

상품을 거래하는 구조에는 \textbf{Revert}가 필요하며, Revert를 암호학적으로
해결하는 과정에서 연구의 기여를 만들 수 있다고 생각했다.

\section{Revert 연구}

\subsection{Revert의 의미와 필요성}

Revert는 \textbf{이미 판매한 상품을 판매자가 직접 회수하는 것}을 의미한다.

\[
  \boxed{\text{Seller A}}
  \;\longrightarrow\;
  \boxed{\text{Buyer B}}
  \qquad
  \overset{\text{Revert}}{\longrightarrow}
  \qquad
  \boxed{\text{Seller A}}
\]

다음과 같은 상황에서 A가 다른 주체에 의존하지 않고, \textbf{A 자신의
  능력으로} 상품을 회수할 수 있어야 한다고 생각했다.

\begin{itemize}
  \item B가 상품 대금을 지불하지 않은 경우
  \item A가 B가 아닌 다른 사용자 \(B'\)에게 상품을 잘못 전송한 경우
  \item 판매 상품의 탄소 배출량, 재활용률 또는 수량을 잘못 입력한 경우
\end{itemize}

\subsection{공급망의 상류와 하류}

상품이 A에서 B를 거쳐 C에게 판매되는 공급망을 다음과 같이 구분했다.

\[
  \boxed{\text{A: upstream}}
  \;\longrightarrow\;
  \boxed{\text{B: A의 direct downstream}}
  \;\longrightarrow\;
  \boxed{\text{C: A의 indirect downstream}}
\]

A는 직접 하류인 B를 알지만, 간접 하류인 C는 알지 못하는 구조로 보았다.

\subsection{간접 하류에서 발생하는 한계}

상품이 C까지 전달되면 A가 해당 상품의 정보를 직접 수정하기 어렵다. 특히 탄소
배출량이나 제품의수량처럼 직접 하류와 간접 하류에 모두 영향을 주는 정보는 상류에서
단순히 변경할 수 없다고 보았다.\\

상품의 부가 정보인 \texttt{DocumentInfo}만 수정하는 경우에도 A는 직접 하류인
B에게 변경 사실을 알릴 수 있지만, 간접 하류인 C에게는 변경된 정보를 전달하기
어렵다. (Pointer-based redact 논문을 응용할 수 있다고 보았음.)\\

정보 수정뿐 아니라 회수를 허용할지 결정하는 과정이 필요하다고도 생각하여 Voting-based redactable 방식도 살펴봤으나
논문의 주제를 크게 벗어난다고 생각하고, 적절하지 않다고 판단했다.

\subsection{결론}

간접하류까지 진행된 공급망을 되돌리는 것은 어렵다고 보았다.


\subsection{현재 프로토콜에서 Revert를 구현한 방법}

A가 직접 하류인 B와 거래한 상품을 다시 회수할 수 있도록, 현재 프로토콜은
\texttt{RecallVoucher}를 사용한다.

\[
  \boxed{\text{A: sender}}
  \;\xrightarrow{\text{Transfer + RecallVoucher}}\;
  \boxed{\text{B: direct downstream}}
\]

A는 \texttt{RecallVoucher}가 제품 수령이나 Recall에 아직 사용되지 않았고
deadline 이전인 경우에 자신의 권한으로 Recall할 수 있다. \(\Delta t\)는
Transfer 시점부터 Recall을 허용하는 \textbf{최대 시간 구간의 길이}다.

\[
  t_{\mathrm{Recall}}
  \;<\;
  t_{\mathrm{Transfer}}+\Delta t
\]

\(\Delta t\)가 지나면 판매자의 Recall 권한은 만료된다.

\subsection{RecallVoucher의 한계}

Transfer
이후 상품은 pending 상태에 놓이고, Buyer의 제품 수령 또는 Seller의 Recall 중 먼저
실행된 동작에 따라 최종 상태가 결정된다.

\textbf{Buyer가 먼저 수령하는 경우.}
Buyer가 \(\Delta t\) 안에 제품 수령을 진행하면
\texttt{RecallVoucher}를 사용한 거래가 확정되고 Seller의 Recall 권한은 사라진다.
따라서 현재 프로토콜은 Seller가 \(\Delta t\) 전체 기간 동안 Recall 권한을
유지하도록 강제하지 않는다.

\[
  \boxed{\text{RecallVoucher: Pending}}
  \;\longrightarrow\;
  \begin{cases}
    \text{Buyer가 제품 수령}   & \Rightarrow \text{Seller의 Recall 권한 소멸} \\
    \text{Seller가 Recall} & \Rightarrow \text{상품 회수}
  \end{cases}
\]

\textbf{Freeze를 적용하는 경우.}
\(\Delta t\) 동안 \(cm\)을 거래하지 못하게 freeze하면 Recall 가능 기간을 보장할
수 있다. 그러나 그 기간 동안 B가 C에게 상품을 전달하지 못하므로, 공급망의 흐름이
\(\Delta t\)만큼 중단되는 문제가 발생한다.

\[
  \boxed{\text{B가 상품 수령}}
  \;\xrightarrow{\text{freeze for }\Delta t}\;
  \boxed{\text{C에게 전달 불가}}
\]

\subsection{참고한 연구}

\textbf{ERC-20R and ERC-721R}은 짧은 이의 제기 기간(dispute window) 동안
분쟁 대상이 된 token transfer를 동결하고, distributed judge의 결정에 따라 해당 거래를 되돌릴 수 있도록 설계한
가역적인 token standard 이다.

현재 프로토콜의 RecallVoucher는 이 논문에서 \textbf{일정 기간 동안만 거래의 회수를 허용하는 개념}을 참고했다.

다만 ERC-20R/721R은 \textbf{judge}의 판단을 거쳐 거래를 되돌리는 반면, 현재 프로토콜에서는
sender가 \(\Delta t\) 안에 자신의 직접 하류와 이루어진 거래를 직접 Recall한다.


\section{현재 프로토콜의 구조}

\subsection{DPP의 등장 배경, 정의와 목적}

제품의 지속가능성 및 생애주기 정보가 제조사, 공급사, 재활용업체 등 여러 조직에 분산되어 있어,
하류 참여자가 제품의 원산지·구성·사용·수리 정보를 신뢰할 수 있는 형태로 확인하기 어렵다. \\

이를 해결하기 위해서 각 물품에 대한 정보를 담는 객체, 즉 디지털 제품 여권인 DPP(Digital Product Passport)가 필요하다는 요구가 제기되었다. \\\\\\\\

EU에서는 DPP를 다음과 같이 설명한다:

\begin{itemize}
  \item DPP는 제품별로 규제를 만족하기 위한 정보를 포함하고, data carrier(QR code 등)를 통해
        전자적으로 접근할 수 있는 특정 제품에 관한 데이터 집합이다.
\end{itemize}

본 논문에서는 DPP를 제품의 생애주기 동안 갱신할 수 있는 동적인 데이터 객체라 정의한다.  \\

DPP는 고유 식별자와 연결되며 제품에 대한 정보(State)는 생애주기 동안 갱신될 수 있지만, 동일한 식별 대상을 나타내는 동안 해당 식별자는 유지된다.

\begin{center}
  \renewcommand{\arraystretch}{1.35}
  \setlength{\tabcolsep}{8pt}
  \begin{tabularx}{\textwidth}{
      @{}
      >{\raggedright\arraybackslash}p{28mm}
      >{\raggedright\arraybackslash}X
      >{\raggedright\arraybackslash}X
      @{}
    }
    \toprule
    \textbf{DPP 식별 수준} & \textbf{식별자 설명}       & \textbf{식별 대상} \\
    \midrule
    Model level
                       & 제품 모델에 부여된 고유 식별자
                       & 동일한 기본 설계를 공유하는 제품 모델                  \\
    Batch/Lot level
                       & 생산 배치에 부여된 고유 식별자
                       & 동일한 생산 배치에 속하는 제품 집합                   \\
    Item level
                       & 개별 제품에 부여된 고유 식별자
                       & 하나의 물리적인 고유한 제품                        \\
    \bottomrule
  \end{tabularx}
\end{center}

따라서 고유 식별자는 항상 개별 제품 하나를 의미하지 않는다.

\subsection{zkDPP의 목표}

zkDPP는 회사 및 제품에 대한 정보를 공개하지 않으면서, 공급망 Event를 거치면서 갱신된 DPP 정보가
이전 정보로부터 올바르게 도출되었음을 검증하는 프로토콜이다.

\subsection{핵심 가정}

1. \textbf{State 정의 가정}

제품 또는 산업별로 관리할 State의 종류, 단위가 사전에 정의되어 있다고 가정한다. \\
정의된 State에 따라서, Merge, Split, Transfer에서 발생할 State Transition이 결정되어 있다고 가정한다. \\

2. \textbf{Entry 신뢰 가정}

TTP가 실제 물품과 최초 DPP 정보의 일치 여부를 확인하거나, 인증된 업체에 Entry 권한을 부여한다. \\
zkDPP는 Entry 이후 발생하는 State transition의 일관성을 검증한다. \\

3. \textbf{Process TTP 신뢰 가정}

TTP가 Process 단계에서 사용할 \(ek, vk\)를 제공한다고 가정한다. \\
Process 단계에서는 입력/출력 제품의 종류, quantity, state가 올바른지 확인한다.\\

4. \textbf{Provenance 추적 가정}

Input과 output commitment의 관계를 추적할 수 있게 만드는 event를 발행하지 않는다.\\

문제가 발생한 \(cm\)을 감사할 때에는 관련 참여자가 보관한 note와 이전 거래 정보를
감사자에게 제공하고, 감사자는 참여자의 협조를 받아 상류 방향으로 이 과정을 반복한다고
가정한다.

\subsection{데이터 모델}

\subsubsection*{1. \(Document\)}

\[
  \mathit{Document}
  =
  (\mathit{DocumentHash},q,\mathbf{a},\mathbf{s})
\]

\begin{itemize}
  \item \(\mathit{DocumentHash}=H(\mathit{DocumentInfo})\):
        Event에서 직접 검증하지 않는 제품 메타데이터의 hash

        \begin{itemize}
          \item \(\mathit{DocumentInfo}=(d_1,d_2,\ldots,d_n)\):
                제품을 설명하기 위한 정보
        \end{itemize}


  \item \(q\):
        물품의 quantity

  \item \(\mathbf{a}=(a_1,a_2,\ldots,a_u)\):
        비산술적 속성 전이를 결정하는 policy attributes

        \begin{itemize}
          \item 예:
                \(\mathit{ProductIdentity}\), \(\mathit{OriginId}\),
                \(\mathit{UnitId}\)
                \begin{itemize}
                  \item \(\mathit{ProductIdentity}\): \{Model, Batch, Item\}
                \end{itemize}
        \end{itemize}

  \item \(\mathbf{s}=(s_1,s_2,\ldots,s_{\ell})\):
        공급망 Event에서 보존, 분배, 증가 또는 감소 관계를 검증하는
        quantitative flow states

        \begin{itemize}
          \item 예:
                \(\mathit{total}_{kg}\), \(\mathit{carbon}_{kgCO2e}\),
                \(\mathit{recycled}_{kg}\)
        \end{itemize}
\end{itemize}

\subsubsection*{2. \(cm\)}

\[
  cm
  =
  H(
  \mathit{Document},
  addr,
  o_{cm}
  )
\]

\[
  sk_{own}\xleftarrow{\$}\mathbb{F}_r,
  \qquad
  pk_{own}=H(sk_{own}),
  \qquad
  addr=H(pk_{own})
\]



\begin{itemize}
  \item \(o_{cm}\):
        randomness

  \item \(addr\):
        \(pk_{own}\)으로부터 도출되는 holder의 address

        \begin{itemize}
          \item \(sk_{own}\):
                holder의 secret key

          \item \(pk_{own}\):
                \(sk_{own}\)으로부터 도출되는 public key
        \end{itemize}
\end{itemize}

\subsubsection*{2.1 \(nf\)}

\[
  nf=H(sk_{own},cm)
\]

\(cm\)의 중복 소비를 방지하는 nullifier

\subsubsection*{2.2 \(note\)}

\[
  note=(cm,\mathit{Document},o_{cm})
\]

Holder가 지갑에 보관하는 note

\subsubsection*{3. \(rv\)}


\[
  rv
  =
  H(
  \mathit{Document},
  addr_{Sender},
  addr_{Receiver},
  deadline_{epoch},
  o_{rv}
  )
\]

\begin{itemize}
  \item \(deadline_{epoch}=transfer_{epoch}+\Delta_{epoch}\):
        Sender가 Recall할 수 있는 마지막 경계를 나타내는 private deadline
        \begin{itemize}
          \item \(current_{epoch}<deadline_{epoch}\):
                Recall 허용 조건
        \end{itemize}
\end{itemize}

Transfer된 \(\mathit{Document}\)가 Proceed 또는 Recall되기 전의 대기 상태를 나타내는 commitment

\subsubsection*{3.1 \(nf_{rv}\)}

\[
  nf_{rv}=H(o_{rv},rv)
\]

Proceed와 Recall 중 한 번만 처리되도록 하는 nullifier

\subsubsection*{3.2 \(note_{rv}\)}

\[
  note_{rv}
  =
  (
  rv,
  \mathit{Document},
  addr_{Sender},
  addr_{Receiver},
  deadline_{epoch},
  o_{rv}
  )
\]

Sender와 Receiver가 동일한 구조로 note를 보관한다.

\subsection{프로토콜 개요}

\begin{center}
  \begin{tabular}{ll}
    \hline
    Event    & Protocol                                          \\
    \hline
    Entry    & $\varnothing\rightarrow cm$                       \\
    Transfer & $cm\rightarrow (rv, cm_{change},)$                \\
    Proceed  & $rv\rightarrow cm$                                \\
    Recall   & $rv\rightarrow cm$                                \\
    Merge    & $(cm_{1},cm_{2})\rightarrow cm$                   \\
    Split    & $cm\rightarrow (cm_{1},cm_{2})$                   \\
    Process  & $(cm_{i})_{i=1}^{m}\rightarrow(cm_{j})_{j=1}^{n}$ \\
    Exit     & $cm\rightarrow\varnothing$                        \\
    \hline
  \end{tabular}
\end{center}

\subsection{Quantity, Policy Attributes, State}

\begin{center}
  \renewcommand{\arraystretch}{1.35}
  \setlength{\tabcolsep}{8pt}
  \begin{tabularx}{\textwidth}{
      @{}
      >{\raggedright\arraybackslash}p{28mm}
      >{\raggedright\arraybackslash}X
      @{}
    }
    \toprule
    \textbf{구성} & \textbf{구분 기준}                    \\
    \midrule
    \(\mathit{DocumentHash}\)
                & Event에서 직접 일관성을 검사하지 않는 제품 정보     \\
    \(q\)
                & 물품의 quantity                      \\
    \(\mathbf{a}\)
                & 제품의 식별단위, 원산지, 수치단위처럼 물품을 구분하는 속성 \\
    \(\mathbf{s}\)
                & 수치적 보존, 분배, 증가 또는 감소 관계를 갖는 State \\
    \bottomrule
  \end{tabularx}
\end{center}

\subsubsection*{Merge}

Merge 가능 여부와 출력 policy attributes는 산업별
\(\mathsf{MergePolicy}\)로 추상화한다.

\[
  \mathsf{MergePolicy}
  (\mathbf{a}_1,\mathbf{a}_2,\mathbf{a}_{out})
  =1
\]

\[
  q_{out}=q_1+q_2,
  \qquad
  \mathbf{s}_{out}=\mathbf{s}_1+\mathbf{s}_2
\]

\subsubsection*{Split}

Split 가능 여부와 출력 policy attributes는
\(\mathsf{SplitPolicy}\)로 추상화한다.

\[
  \mathsf{SplitPolicy}
  (\mathbf{a}_{in},\mathbf{a}_{out,1},\mathbf{a}_{out,2})
  =1
\]

\[
  q_{in}=q_1+q_2,
  \qquad
  \mathbf{s}_{in}=\mathbf{s}_1+\mathbf{s}_2
\]

\(\mathsf{MergePolicy}\)와 \(\mathsf{SplitPolicy}\)는 산업군별로 결정되며,
policy를 검증하는 고정된 vk를 사용한다.

\subsubsection*{Process}

Process는 입력 물품을 소비하고 새로운 물품을 생성할 수 있으므로
\(q\), \(\mathbf{a}\), \(\mathbf{s}\)의 전이를 함께 검증한다.

\[
  \mathsf{ProcessPolicy}
  (
  \mathbf{q}_{in},
  \mathbf{a}_{in},
  \mathbf{s}_{in},
  \mathbf{q}_{out},
  \mathbf{a}_{out},
  \mathbf{s}_{out}
  )
  =1
\]

Process는 공정마다 서로 다른 입력과 출력 조건을 가질 수 있다.
TTP는 공정별 \(\mathsf{ProcessPolicy}\)를 검증하는 vk를 매번 등록한다.

\subsection{8개 Event의 동작}

\subsubsection*{Entry}

\[
  \varnothing
  \xrightarrow{\mathrm{Entry}}
  cm_{new}
\]

\begin{enumerate}[label=\textbf{\arabic*.},leftmargin=18pt]
  \item Issuer TTP 또는 인증된 업체가 실제 물품과 최초
        \(\mathit{Document}\)의 일치 여부를 확인한다.
  \item 최초 holder가 \((\mathit{Document},addr,o_{cm})\)로
        \(cm_{new}\)를 생성한다.
  \item Contract는 중복되지 않은 \(cm_{new}\)를 \(MT\)에 등록한다.
\end{enumerate}

\subsubsection*{Transfer}

\[
  cm_{old}
  \xrightarrow{\mathrm{Transfer}}
  (rv_{new},cm_{change})
\]

\begin{enumerate}[label=\textbf{\arabic*.},leftmargin=18pt]
  \item Sender는 자신이 소유한 \(cm_{old}\)에서 Receiver에게 전달할 quantity와
        자신에게 남길 change quantity를 정한다.
  \item \(\mathbf{a}\)는 유지하고, \(\mathbf{s}\)는 두 quantity에 따라
        transfer 부분과 change 부분으로 분배한다.
  \item proof는 \(MT\) membership, nf, quantity와
        State 분배, \(rv_{new}\)와 \(cm_{change}\)의 생성을 검증한다.
  \item Contract는 \(cm_{old}\)를 소비하고, \(cm_{change}\)를 \(MT\)에,
        \(rv_{new}\)를 \(rvMT\)에 등록한다.
\end{enumerate}

\subsubsection*{Proceed}

\[
  rv_{old}
  \xrightarrow{\mathrm{Proceed}}
  cm_{receiver}
\]

\begin{enumerate}[label=\textbf{\arabic*.},leftmargin=18pt]
  \item Receiver는 전달받은 \(note_{rv}\)와 \(rvMT\) membership path를 준비한다.
  \item Receiver는 \(rv_{old}\)로부터 \(nf_{rv}\)를 만들고, 동일한
        \(\mathit{Document}\)를 자신에게 귀속하는 \(cm_{receiver}\)를 생성한다.
  \item Proceed proof는 Receiver의 권한, \(rv_{old}\)의 membership,
        \(nf_{rv}\)와 \(cm_{receiver}\)의 생성을 검증한다.
  \item Contract는 \(nf_{rv}\)를 등록하여 해당 voucher의 재사용을 막고,
        \(cm_{receiver}\)를 \(MT\)에 등록한다.
\end{enumerate}

\subsubsection*{Recall}

\[
  rv_{old}
  \xrightarrow{\mathrm{Recall}}
  cm_{sender}
\]

\begin{enumerate}[label=\textbf{\arabic*.},leftmargin=18pt]
  \item Sender는 자신이 생성했던 \(note_{rv}\)와 \(rvMT\) membership path를
        준비한다.
  \item Sender는 \(rv_{old}\)로부터 \(nf_{rv}\)를 만들고, 동일한
        \(\mathit{Document}\)를 자신에게 되돌리는 \(cm_{sender}\)를 생성한다.
  \item Recall proof는 Sender의 권한, \(rv_{old}\)의 membership,
        \(current_{epoch}<deadline_{epoch}\), \(nf_{rv}\)와
        \(cm_{sender}\)의 생성을 검증한다.
  \item Contract는 \(nf_{rv}\)를 등록하여 Proceed와 Recall의 중복 실행을 막고,
        \(cm_{sender}\)를 \(MT\)에 등록한다.
\end{enumerate}

\subsubsection*{Merge}

\[
  (cm_{old,1},cm_{old,2})
  \xrightarrow{\mathrm{Merge}}
  cm_{new}
\]

\begin{enumerate}[label=\textbf{\arabic*.},leftmargin=18pt]
  \item Holder는 자신이 소유한 서로 다른 두 finalized note를 선택한다.
  \item 배포 시 고정된 \(\mathsf{MergePolicy}\)가
        \((\mathbf{a}_1,\mathbf{a}_2,\mathbf{a}_{new})\)의 조합을 허용하는지
        검사한다.
  \item Merge proof는 두 input의 소유권과 membership,
        \(q_{new}=q_1+q_2\), \(\mathbf{s}_{new}=\mathbf{s}_1+\mathbf{s}_2\),
        \(cm_{new}\)의 생성을 검증한다.
  \item Contract는 두 input nullifier를 등록하고 \(cm_{new}\)를 \(MT\)에
        등록한다.
\end{enumerate}

\subsubsection*{Split}

\[
  cm_{old}
  \xrightarrow{\mathrm{Split}}
  (cm_{new,1},cm_{new,2})
\]

\begin{enumerate}[label=\textbf{\arabic*.},leftmargin=18pt]
  \item Holder는 input quantity를 두 output quantity로 나눈다.
  \item 배포 시 고정된 \(\mathsf{SplitPolicy}\)가
        \((\mathbf{a}_{old},\mathbf{a}_{new,1},\mathbf{a}_{new,2})\)의
        전이를 허용하는지 검사한다. 동일한 Item-level identifier를 두 output에
        복제하는 Split은 허용하지 않는다.
  \item Split proof는 소유권과 membership, quantity 보존,
        \(\mathbf{s}\)의 분배와 residual 처리, 두 output commitment의 생성을
        검증한다.
  \item Contract는 input nullifier를 등록하고 두 output commitment를 순서대로
        \(MT\)에 등록한다.
\end{enumerate}

\subsubsection*{Process}

\[
  (cm_{old,i})_{i=1}^{m}
  \xrightarrow{\mathrm{ProcessPolicyId}}
  (cm_{new,j})_{j=1}^{n}
\]

\begin{enumerate}[label=\textbf{\arabic*.},leftmargin=18pt]
  \item Processor는 사용할 공정의 \(\mathit{ProcessPolicyId}\), input note와
        생성할 output을 선택한다.
  \item Contract는 \(\mathit{ProcessPolicyId}\)에 등록된 Process verifier를
        선택한다.
  \item Process proof는 input의 소유권과 membership, 중복 사용 방지,
        \((\mathbf{q}_{in},\mathbf{a}_{in},\mathbf{s}_{in})\)에서
        \((\mathbf{q}_{out},\mathbf{a}_{out},\mathbf{s}_{out})\)으로의 전이가
        해당 공정 policy를 만족하는지 검증한다.
  \item Contract는 모든 input nullifier를 등록하고, 생성된 output commitment를
        순서대로 \(MT\)에 등록한다.
\end{enumerate}

\subsubsection*{Exit}

\[
  cm_{old}
  \xrightarrow{\mathrm{Exit}}
  \varnothing
\]

\begin{enumerate}[label=\textbf{\arabic*.},leftmargin=18pt]
  \item Holder는 공급망에서 제거할 finalized note를 선택한다.
  \item Exit proof는 해당 note의 소유권, \(MT\) membership과 올바른
        nullifier 생성을 검증한다.
  \item Contract는 nullifier가 아직 사용되지 않았는지 확인하고 proof를 검증한다.
  \item Contract는 nullifier만 등록하며 새로운 commitment는 생성하지 않는다.
\end{enumerate}

\subsection{고민 사항}

\end{document}
